\documentclass[12pt,UTF-8]{UtopiaPhyCyclus_book}

% ==================================================
% 答案显示模式设置
% 0: 正常显示答案 (Show Answer)
% 1: 隐藏答案但保留篇幅空白 (Phantom Space)
% 2: 隐藏答案并保留固定空白 (Fixed Space)
% 3: 完全隐藏答案不留白 (No Space)
\def\solutionmode{0} %在此处修改数字即可切换模式

\NewDocumentEnvironment{solutions}{ O{【答案】} O{} O{} +b }{%
  \ifcase\solutionmode\relax
    % Mode 0: 正常显示
    \par\noindent{\textcolor{main}{\bfseries \if\relax\detokenize{#1}\relax 【答案】\else #1\fi}}\,%
    #4%
    \par\vspace{\baselineskip}%
  \or
    % Mode 1: 留白篇幅长度
    \par\noindent{\textcolor{main}{\bfseries \if\relax\detokenize{#1}\relax 【答案】\else #1\fi}}\,%
    \phantom{%
      \parbox[t]{\linewidth}{%
         \renewenvironment{figure}[1][]{\center\captionsetup{type=figure}}{\endcenter}%
         \renewenvironment{table}[1][]{\center\captionsetup{type=table}}{\endcenter}%
         #4%
      }%
    }%
    \par\vspace{\baselineskip}%
  \or
    % Mode 2: 留白固定长度 (默认 5cm，可按需调整)
    \par\noindent{\textcolor{main}{\bfseries \if\relax\detokenize{#1}\relax 【答案】\else #1\fi}}\,%
    \vspace*{5cm}%
    \par\vspace{\baselineskip}%
  \or
    % Mode 3: 不留白 (什么都不显示)
    % Do nothing
  \fi
}{}

\NewDocumentEnvironment{answer}{ O{【答案】} O{} O{} +b }{%
  \ifcase\solutionmode\relax
    % Mode 0: 正常显示
    \par\noindent{\textcolor{main}{\bfseries \if\relax\detokenize{#1}\relax 【答案】\else #1\fi}}\,%
    #4%
    \par\vspace{\baselineskip}%
  \or
    % Mode 1: 不留白 (什么都不显示)
  \fi
}{}

% ==================================================


% ==================================================
% 自动替换 \left( \right) 为 ( ) 和 \left[ \right] 为 [ ]
% 保持其他 delimiter 的 \left \right 行为不变
% ==================================================
\let\oldleft\left
\let\oldright\right
\def\left#1{%
  \ifx#1(\mathopen{(}\else
  \ifx#1[\mathopen{[}\else
  \oldleft#1\fi\fi}
\def\right#1{%
  \ifx#1)\mathclose{)}\else
  \ifx#1]\mathclose{]}\else
  \oldright#1\fi\fi}

% 标题设置

\title{高中物理思维体系}
\author{槿灵兮}
\date{\today}

\begin{document}



\thispagestyle{nohf}
\includepdf[pages=1,fitpaper=true]{cover/cover.pdf}

%\cleardoublepage  % 确保下一页从奇数页开始
% 添加超链接选项以避免PDF字符串中的令牌错误


\thispagestyle{nohf}
\clearpage
\tableofcontents
\thispagestyle{nohf}
\newpage

\setcounter{page}{1}


\chapter{运动的描述}
\input{chapters/运动的描述/content_运动的描述}

\chapter{匀变速直线运动}
\input{chapters/匀变速直线运动/content_匀变速直线运动}

\chapter{平衡力学}
\input{chapters/平衡力学/content_平衡力学}

\chapter{牛顿运动定律}
\input{chapters/牛顿运动定律/content_牛顿运动定律}

\chapter{曲线运动}
\input{chapters/曲线运动/content_曲线运动}

\chapter{万有引力与航天}
\input{chapters/万有引力与航天/content_万有引力与航天}

\chapter{功与能量}
\input{chapters/功与能量/content_功与能量}

\chapter{动量}
\input{chapters/动量/content_动量}

\chapter{电场}
\input{chapters/电场/content_电场}

\chapter{电路}
\input{chapters/电路/content_电路}

\chapter{磁场运动}
\input{chapters/磁场运动/content_磁场运动}

\chapter{电磁感应}
\input{chapters/电磁感应/content_电磁感应}

\chapter{交流电}
\input{chapters/交流电/content_交流电}

\chapter{机械振动与机械波}
\input{chapters/机械振动与机械波/content_机械振动与机械波}

\chapter{光学}
\input{chapters/光学/content_光学}

\chapter{热学}
\input{chapters/热学/content_热学}

\chapter{近代物理}
\input{chapters/近代物理/content_近代物理}

\chapter{电学实验}
\input{chapters/电学实验/content_电学实验}

\chapter{力学试验}
\input{chapters/力学试验/content_力学试验}

\chapter{选修实验}
\input{chapters/选修实验/content_选修实验}

\end{document}
